\mode*
\section{The remote shell: SSH}
\label{sec:ssh}

The fourth educational objective is the remote shell: that a shell can just
as well be attached to another machine, and what follows from that.

\subsection{Documented difficulties}

The seed literature round found \emph{no} studies of how learners
(mis)understand SSH or remote shells---a genuine gap.
% TODO: The systematic phase must confirm or fill this gap (queries on
% remote access, SSH education, distributed-systems misconceptions at the
% tool level).
What we have are recurring observations from our own course and from the
university's IT department: students lose track of which machine their
shell is attached to, and chain \texttt{ssh} sessions into cycles
(\texttt{ssh A}, then \texttt{ssh B}, then \texttt{ssh A} again) instead
of opening parallel sessions.
% TODO: Collect and cite these observations properly (course issue
% tracker; IT correspondence), clearly marked as course data rather than
% literature.

\subsection{Candidate critical aspects}

% XXX Preliminary.
\begin{itemize}
  \item \emph{A shell runs somewhere}: every prompt belongs to exactly one
    machine; the prompt (and \texttt{hostname}) tells you where you are.
  \item \emph{Same shell, different computer}: the interface is invariant,
    the machine varies---which is precisely why it is easy to lose track.
  \item \emph{Files live on machines}: a file created on the remote host
    is not on the laptop, and vice versa.
\end{itemize}

\subsection{Patterns of variation}

\begin{itemize}
  \item Contrast (machine varies, shell invariant): run
    \texttt{hostname}, \texttt{ls} before and after \texttt{ssh}; same
    commands, different answers.
  \item Contrast (session structure varies): a chained
    \texttt{ssh A -> ssh B} versus two parallel terminals; observe
    \texttt{exit} behaviour.
  \item Fusion: edit locally, copy to the remote host, run remotely---the
    full round trip.
\end{itemize}
